\documentclass[11pt,a4paper]{article}
\usepackage[utf8]{inputenc}
\usepackage[T1]{fontenc}
\usepackage{amsmath,amssymb,amsthm}

\usepackage{graphicx}
\usepackage{booktabs}
\usepackage{array}
\usepackage{xcolor}
\usepackage{hyperref}
\hypersetup{colorlinks=true,linkcolor=blue!50!black,citecolor=blue!50!black,urlcolor=blue!50!black}
\usepackage[a4paper,margin=1in]{geometry}


\usepackage{enumitem}
\usepackage{caption}
\usepackage{subcaption}
\usepackage{titlesec}
\titlespacing\section{0pt}{14pt plus 4pt minus 2pt}{6pt plus 2pt minus 2pt}
\titlespacing\subsection{0pt}{10pt plus 4pt minus 2pt}{4pt plus 2pt minus 2pt}

% Math shortcuts
\newcommand{\R}{\mathbb{R}}
\newcommand{\E}{\mathbb{E}}
\newcommand{\norm}[1]{\left\|#1\right\|}
\newcommand{\inner}[2]{\left\langle#1,#2\right\rangle}
\newcommand{\Ins}{\text{Ins}}
\newcommand{\Ret}{\text{Ret}}
\newcommand{\Com}{\text{Com}}

\title{\bf Addressable Memory: A Lossless Primitive for \\
Sequential Knowledge Editing in Frozen Language Models}

\author{Mohamed Mabrok \\ Minds-R-Lab}

\date{\today}

\begin{document}
\maketitle

\begin{abstract}
\noindent
Adding new factual knowledge to a deployed language model is the most basic
form of learning a database can do --- and yet every weight-modification
method we know (sequential fine-tuning, LoRA, and rank-one knowledge editors
like ROME / MEMIT) destroys prior knowledge as the number of edits grows.
We argue this is not a tuning problem but a structural one: facts written
into shared parameters necessarily compete with the parameters' existing
function. We propose an alternative primitive --- an \emph{addressable
memory layer} --- that wraps a single mid-stack MLP and stores each new fact
as a $(k, \delta v)$ pair in an external bank, leaving every base weight
untouched. The key $k$ is the MLP intermediate at the last prompt token
captured during a teacher-forced insertion pass; the value $\delta v$ is a
gradient-optimized perturbation that drives the language modeling head to
the target. At inference, cosine similarity over keys gates which slot
fires per position. We introduce SFIB, a pre-registered benchmark for
sequential one-shot fact insertion over a synthetic knowledge base of
2{,}000 pretrain triples and 500 insertion triples with fictional
syllable-pool entities. On GPT-2 small (124M), with 3-seed validation, our
primitive achieves Insertion@500 = $0.468 \pm 0.030$ and Retention@500 =
$0.922 \pm 0.033$ --- $3.2\times$ MEMIT on insertion and $5.9\times$ on
retention, simultaneously, over 500 sequential edits. We further extend
the primitive with multi-position $v^\star$-optimization to handle
multi-token target answers under modern BPE tokenizers, and demonstrate
on Qwen2.5-0.5B that the architecture transfers to a different family
(SwiGLU MLP, RoPE, GQA) achieving Insertion@500 = $0.995$ at
Retention@500 = $0.811$ --- the only weight-modification method we tested
that preserves more than $5\%$ of prior knowledge under this stream.
\end{abstract}

\section{Introduction}

A language model that cannot learn a new fact after deployment is, in a
narrow but important sense, less capable than a database. Databases insert
records in $O(1)$ time without disturbing other records. Language models
require gradient descent through tens to hundreds of billions of parameters
to absorb a single new fact, and the gradient invariably disturbs
neighboring memories. This trade-off has been the central technical
obstacle to continual learning for as long as neural networks have existed,
and it has shaped a generation of workarounds --- replay buffers, parameter
freezing, low-rank adapters, retrieval-augmented generation, and explicit
knowledge editors that try to make minimal weight changes per fact.

The state-of-the-art editor at the time of writing is the MEMIT family of
methods \citep{meng2022rome,meng2022memit}, which uses causal tracing to
localize a single fact to a particular mid-stack MLP and then applies a
rank-one update to that MLP's down-projection. When a single fact is
inserted in isolation, MEMIT can achieve 99\% efficacy with negligible
collateral. The original paper reports successful editing of thousands of
facts in one batched update. Yet \emph{sequential} insertion --- the
realistic deployed setting where facts arrive one at a time over the life
of a model --- shows a different picture. Each rank-one update overlaps
in support with all previous updates, the perturbations accumulate in the
same low-dimensional subspaces of MLP weight space, and accuracy on both
the new facts and on the old pretrained knowledge degrades together.

We argue that this degradation is fundamental, not technical. Any method
that writes new information into a frozen-but-overwritable parameter set
must spend that parameter set's degrees of freedom; eventually the
parameters cannot simultaneously represent everything they used to plus
everything they are now asked to know. The only way around this is to
\emph{not write into the existing parameters at all} --- to instead
maintain new information in a separately-addressable structure that the
model can consult without disturbing its prior weights.

This idea is not new in form. Memory-augmented neural networks date back
at least to the Neural Turing Machine \citep{graves2014ntm} and the
Differentiable Neural Computer \citep{graves2016dnc}; Memory Networks
\citep{weston2015memnets} and Modern Hopfield Networks
\citep{ramsauer2020modern} explore variants of the same theme. Most
recently, \citet{berges2024memory} scaled product-key memory layers to
trillion-token training budgets and showed they outperform dense FFN
layers on knowledge-intensive evaluation. What has been less explored is
the specific case of \emph{post-training, one-shot, sequential insertion}
on a frozen pretrained backbone, where the question is not whether
addressable memory can help during pretraining but whether it can serve as
the alternative to gradient-based editors like MEMIT after a model is
already in deployment. This is the regime we study.

Our contributions are four:

\begin{enumerate}[leftmargin=*,topsep=2pt,itemsep=2pt]
\item We introduce the \textbf{Sequential Fact Insertion Benchmark
(SFIB)}: a pre-registered synthetic knowledge base with 2{,}000 pretrain
triples, 500 insertion triples (disjoint entities), and 200 multi-hop
composition queries. SFIB measures three quantities at every checkpoint
$N$ along the insertion stream --- $\Ins@N$ (accuracy on the inserted
facts), $\Ret@N$ (accuracy on the pretrained facts), and $\Com@N$
(accuracy on two-hop questions combining one old fact and one new fact).
Entities are constructed from fictional syllable pools to guarantee no
overlap with the model's pretraining distribution.

\item We propose the \textbf{addressable memory primitive}: a drop-in
wrapper around a chosen mid-stack MLP that stores each new fact as a
$(k, \delta v)$ slot in an external key-value bank. At forward pass time,
cosine similarity between the layer's intermediate state and stored keys
selects which slot's $\delta v$ to add to the MLP output. No base weight
is ever modified. The primitive is family-agnostic: we test it on GPT-2's
vanilla MLP and on Qwen2.5's SwiGLU MLP using the same code path.

\item We provide \textbf{architectural ablations} showing the contribution
of each design choice. The most important is the firing rule: when slots
fire at \emph{every} sequence position rather than only at the next-token
prediction position, accuracy is capped because spurious slot firings at
upstream tokens pollute the prediction. Restricting firing to the last
position of each forward pass --- which by construction is where the
language-modeling head reads from --- recovers an order-of-magnitude
improvement in Insertion@N at small $N$.

\item We show the architecture \textbf{transfers across model families
and target tokenization regimes}. With a single-position $v^\star$
optimization, the primitive works perfectly when target answers tokenize
to a single subword (as is common under GPT-2's BPE on our synthetic
nouns). For target answers that BPE-tokenize into multiple subwords
(common under Qwen2's tokenizer), single-position perturbation controls
only the first generated token. We extend the primitive with
\emph{multi-position $v^\star$-optimization}: one $(k, \delta v)$ slot is
captured and optimized per predicting position in the prompt-plus-target
sequence. At eval, the wrapped MLP cosine-matches each generated token's
intermediate state against all stored keys, and the slot for that position
fires deterministically. With this extension Qwen2.5-0.5B achieves
Insertion@500 = $0.995$.
\end{enumerate}

The headline number on GPT-2 small with three-seed validation is
$\Ins@500 = 0.468 \pm 0.030$, $\Ret@500 = 0.922 \pm 0.033$. The strongest
baseline (MEMIT, with the same v$^\star$ optimization infrastructure but
the rank-one weight update kept in place) reaches $\Ins@500 = 0.145 \pm
0.031$, $\Ret@500 = 0.186 \pm 0.073$. The next-strongest
(sequential fine-tuning at the canonical learning rate) reaches $\Ins@500
= 0.118$, $\Ret@500 = 0.079$. Across both axes simultaneously, no
weight-modification baseline tested comes within a factor of three of
the addressable memory primitive.

The remainder of the paper is organized as follows. Section~2
reviews related work in knowledge editing and memory-augmented
architectures. Section~3 reviews the relevant background on how
transformer MLPs store and retrieve factual associations. Section~4
introduces the addressable memory primitive formally and gives the
forward and insertion algorithms. Section~5 describes the SFIB benchmark.
Section~6 lays out the experimental protocol. Section~7 reports the
main results, ablations, and the multi-architecture transfer.
Section~8 discusses limitations and connections to related approaches.
Section~9 concludes.

\section{Related Work}

\paragraph{Locate-and-edit knowledge editors.}
The line of work directly comparable to ours is the localize-and-edit
family: \textbf{ROME} \citep{meng2022rome} performs a single rank-one
update to a chosen MLP's down-projection to inject one fact at a time,
guided by causal tracing of where the model stores subject information.
\textbf{MEMIT} \citep{meng2022memit} extends ROME to batched updates
through a least-squares optimum across many edits at once, scaling to
thousands of facts in a single optimization step. Both rely on a $v^\star$
optimization that finds a target MLP-output vector reproducing the desired
prediction, then expresses the required parameter change in closed form.
We use the same $v^\star$-optimization machinery (Adam steps on a
perturbation vector with $L_2$ regularization and norm constraint) but
\emph{store} the result in an external bank instead of dissolving it into
the weights. Subsequent comparisons \citep{hartvigsen2023grace,
hase2023doesedits} have noted that locate-and-edit methods compound
destructively under sequential application, motivating the present work.

\paragraph{Memory-augmented architectures.}
The Neural Turing Machine \citep{graves2014ntm}, the Differentiable Neural
Computer \citep{graves2016dnc}, and the broader Memory Network
\citep{weston2015memnets} family established that recurrent computation
with explicit addressable storage can solve tasks --- algorithmic copy,
sort, associative recall --- that are difficult for vanilla recurrence.
Modern Hopfield Networks \citep{ramsauer2020modern} showed that
exponentially-large autoassociative storage can be implemented with a
small differentiable energy function. Most recent and most directly
relevant is \textbf{Memory Layers at Scale} \citep{berges2024memory},
which replaces a transformer FFN with a learned product-key memory bank;
at trillion-token training scales, knowledge-intensive evaluation
accuracy improves substantially over the dense baseline. Our work
differs in that we are not training a new model: we wrap a single
pre-existing MLP layer of a frozen pretrained model and use the wrapper
only for sequential post-training fact insertion. The memory bank is
populated by a deterministic insertion procedure rather than by joint
training with the base model.

\paragraph{Knowledge in transformer MLPs.}
\citet{geva2021keyvalue} interprets transformer FFN layers as
implicit key-value memories where the first projection's rows act as
keys and the second projection's columns as values. This framing
motivated ROME's choice of editing the MLP down-projection, and it
motivates our choice of MLP intermediate (the activated input to the
down-projection) as our key representation. The crucial difference from
\citet{geva2021keyvalue} is that we make the key-value memory
\emph{external and explicit} rather than implicit in the weight matrix.

\paragraph{Parameter-efficient fine-tuning.}
LoRA \citep{hu2022lora} and adapter-style methods constrain updates to
low-rank or low-parameter subspaces. They reduce the volume of weight
modification per fact but do not change the fundamental property that
new information is written to shared parameters. We confirm this on SFIB:
LoRA-sequential at tuned hyperparameters reaches
$\Ins@500 = 0.008$ and $\Ret@500 = 0.004$ --- worse than full
sequential fine-tuning under the same conditions. The cause is that
forcing each fact's gradient into a few hundred thousand shared
adapter parameters concentrates the perturbation rather than
distributing it; the catastrophe arrives sooner, not later.

\paragraph{Activation steering and controller families.}
A complementary line of recent work intervenes on the residual stream
rather than the weights. \textbf{NTK-Mirror} \citep{chlon2026ntkmirror}
learns sparse signed log-gates on residual-stream channels and composes
them across "controller memories" for different documents or
task styles. The mechanism is multiplicative rather than additive and
operates per-channel rather than per-position. The author explicitly
scopes the method out of factual insertion --- "for facts, use
prompt/RAG provenance" --- so the present work occupies a complementary
slot in design space: per-fact, additive at one MLP layer, with
deterministic insertion via $v^\star$-optimization, targeting the
factual case the controller family declines.

\section{Background: Knowledge in Transformer MLPs}

A standard pre-norm decoder block of a transformer comprises an
attention sublayer and an MLP sublayer, each wrapped in a residual
connection:
\begin{align}
x_\ell^{(\text{mid})} &= x_\ell + \text{Attn}(\text{LN}(x_\ell)), \\
x_{\ell+1} &= x_\ell^{(\text{mid})} + \text{MLP}(\text{LN}(x_\ell^{(\text{mid})})).
\end{align}
The MLP can take one of two forms depending on the architecture family.
For GPT-2 \citep{radford2019gpt2} and other classical decoders:
\begin{equation}
\text{MLP}_{\text{vanilla}}(x) = c_\text{proj}(\sigma(c_\text{fc}(x))),
\end{equation}
where $c_\text{fc}$ and $c_\text{proj}$ are 1-D convolutions
(equivalent to linear layers up to a transpose), and $\sigma$ is a
nonlinearity (GeLU for GPT-2). For Llama \citep{touvron2023llama}, Qwen2
\citep{yang2024qwen2}, and other modern SwiGLU-family decoders:
\begin{equation}
\text{MLP}_{\text{SwiGLU}}(x) = W_\text{down} \left[ \sigma(W_\text{gate} x) \odot W_\text{up} x \right],
\end{equation}
where $\odot$ is elementwise product and $\sigma$ is typically SiLU.

In both forms, the MLP's output at a given position $t$ depends only on
the hidden state at position $t$, computed by the preceding attention.
We will call the intermediate vector immediately before the final
projection ($\sigma(c_\text{fc}(x))$ for vanilla,
$\sigma(W_\text{gate} x) \odot W_\text{up} x$ for SwiGLU) the
\emph{MLP key} at that position, and write it $k_t \in \R^{d_\text{mlp}}$.

\citet{meng2022rome} introduced the methodology of causal tracing to
identify the layer at which a model's prediction for a factual
completion is most sensitive to the subject token's representation.
For GPT-2 medium they found this to be roughly layer 17 of 24, around
70\% of the way through the stack. We use this proportional rule to
choose layers across models: layer 5 of 12 (42\%) for GPT-2 small and
layer 17 of 24 (71\%) for Qwen2.5-0.5B. We confirm with a layer sweep
that adjacent layer choices give qualitatively similar behavior,
suggesting the architecture is not delicately tuned to any specific
storage layer.

The $v^\star$-optimization technique introduced by ROME computes a
target value $v^\star$ for the MLP output at the subject token position
such that, with $v^\star$ substituted in place of the MLP's natural
output at inference, the model produces the desired factual completion.
Concretely, $v^\star$ is found by gradient descent on a free vector
$\delta v$ injected via a forward hook at the chosen position:
\begin{equation}
\delta v^\star = \arg\min_{\delta v} \mathcal{L}_\text{CE}\left(p_\theta(\text{target} \mid \text{prompt}, +\delta v\,@\,\text{pos}), \text{target}\right) + \lambda \frac{\norm{\delta v}^2}{\norm{v_\text{orig}}^2},
\end{equation}
subject to $\norm{\delta v} \le \gamma \norm{v_\text{orig}}$. Following
ROME we use Adam with $\gamma = 4$ and $\lambda = 0.5$ as defaults for
GPT-2; we will see that the SwiGLU family requires more permissive
constraints. ROME then commits the change by computing a rank-one
update to the down-projection that reproduces the same MLP output
at the same MLP key:
\begin{equation}
\Delta W_\text{down} = \frac{(v^\star - W_\text{down} k - b)\, k^\top}{\norm{k}^2}.
\end{equation}
This update lives in the row space spanned by $k$ in $W_\text{down}$ and
therefore perturbs every other input to $W_\text{down}$ that has nonzero
projection onto $k$. As facts accumulate, the rank-one updates
collectively span an increasing fraction of the down-projection's row
space and inevitably interfere with the original weight matrix's encoding
of pretrained associations. We will quantify this interference empirically
in Section~7.

\section{Method: The Addressable Memory Primitive}

\subsection{Architecture}

The primitive is a wrapper module that replaces the MLP submodule at
a chosen layer $\ell$. Let $M$ denote the original (frozen) MLP. The
wrapped module $\hat M$ maintains an external memory bank:
\begin{align}
K &\in \R^{N_\text{max} \times d_\text{mlp}}, \\
V &\in \R^{N_\text{max} \times d_\text{model}},
\end{align}
along with a slot counter $n \in \{0, 1, \ldots, N_\text{max}\}$. The
forward pass at training input $x \in \R^{B \times T \times d_\text{model}}$
computes:
\begin{enumerate}[leftmargin=*,topsep=2pt,itemsep=1pt]
\item Run the base MLP normally to obtain the intermediate state $h$
(shape $B \times T \times d_\text{mlp}$, our keys for retrieval) and the
base output $y_\text{base} = M(x)$ (shape $B \times T \times d_\text{model}$).
\item If $n > 0$, take only the last sequence position $T-1$ and compute
cosine similarities
\begin{equation}
s_b = \frac{h[b, T-1, :]}{\norm{h[b, T-1, :]}} \cdot \frac{K[:n, :]^\top}{\norm{K[:n, :]}_\text{row}}
\in \R^n,
\end{equation}
for each batch element $b$. Let $i_b^\star = \arg\max_i s_b[i]$ and
$g_b = \mathbb{1}[s_b[i_b^\star] > \tau]$, where $\tau$ is the
similarity threshold (default $0.7$).
\item Construct an addition tensor $C \in \R^{B \times T \times d_\text{model}}$
that is zero everywhere except at position $T-1$, where
$C[b, T-1, :] = g_b \cdot V[i_b^\star, :]$.
\item Return $y_\text{base} + C$, applying the MLP's family-specific
trailing dropout if present.
\end{enumerate}
The base MLP's parameters are never modified. The only learnable
quantities in our system are $K, V, n$, all populated by the insertion
procedure described next.

\subsection{Single-position Insertion}

Given a new fact represented as a triple $(s, r, o)$ (subject,
relation, object), the insertion procedure constructs a rewrite prompt
$p$ and target completion $t$ from a template for relation $r$. For our
benchmark, $p$ is in question-answering form
("Where does Pirnebrand live? Answer:") and $t$ is the object with a
leading space ("\,Hamburg"). Let $P$ be the length of the tokenized
prompt and $T$ the length of the tokenized target; the full input is
the concatenation $x_\text{full} \in \mathbb{Z}^{P+T}$.

The insertion procedure runs three forward passes through the model:

\paragraph{(1) Key-value capture.} Run $\text{model}(x_\text{full})$ with
forward hooks on the wrapped MLP at layer $\ell$. From the down-projection
input at position $P-1$ (the last prompt token), record
$k = h_{\ell, P-1}$. From the wrapped MLP's output at the same position,
record $v_\text{orig} = M(\cdot)_{\ell, P-1}$.

\paragraph{(2) $v^\star$-optimization.} Initialize $\delta v = 0 \in
\R^{d_\text{model}}$ as a leaf tensor with gradients enabled. Register a
forward hook on the wrapped MLP that adds $\delta v$ to the MLP output
at position $P-1$ in each forward pass. Run $n_\text{steps}$ steps of
Adam at learning rate $\alpha$ on the joint loss
\begin{equation}
\mathcal{L}(\delta v) = \mathcal{L}_\text{CE}(x_\text{full}; \theta, \text{hook}(\delta v)) + \lambda \frac{\norm{\delta v}^2}{\norm{v_\text{orig}}^2},
\end{equation}
where the cross-entropy is computed on the target token positions
$P, P+1, \ldots, P+T-1$ (the prompt positions are masked with $-100$).
After each gradient step, project $\delta v$ back inside the ball of
radius $\gamma \norm{v_\text{orig}}$:
\begin{equation}
\delta v \leftarrow \min\left(1, \frac{\gamma \norm{v_\text{orig}}}{\norm{\delta v}}\right) \cdot \delta v.
\end{equation}

\paragraph{(3) Slot commit.} Append $k$ and $\delta v$ to row $n$ of the
memory bank and increment $n$.

\subsection{Multi-position Extension}

Single-position insertion works whenever the target $t$ tokenizes to a
single subword: the $\delta v$ at position $P-1$ controls the logit
distribution from which the first (and only) target token is drawn.
When $t$ tokenizes to $T > 1$ subwords, single-position $\delta v$ can
only control the first generated token. Subsequent generated tokens are
predicted from positions $P, P+1, \ldots, P+T-2$, whose hidden states we
have not perturbed. We observe empirically (Section~7) that this
limitation manifests as correct-first-token, incorrect-rest behavior on
multi-token targets under modern BPE tokenizers.

The multi-position extension stores one slot per predicting position.
Let $\mathcal{P} = \{P-1, P, \ldots, P+T-2\}$ index the positions whose
hidden states are responsible for predicting each successive target
token. In step (1) we capture
$\{k_p = h_{\ell, p}, v_{\text{orig}, p} = M(\cdot)_{\ell, p}\}_{p \in \mathcal{P}}$.
In step (2) we maintain $T$ independent perturbation vectors
$\{\delta v_p\}_{p \in \mathcal{P}}$ and add each $\delta v_p$ at position
$p$ via the same forward hook. The CE loss is computed jointly on all
target positions, and the regularizer becomes
$\sum_{p \in \mathcal{P}} \lambda \norm{\delta v_p}^2 / \norm{v_{\text{orig}, p}}^2$.
In step (3) we commit $T$ slots, one per position in $\mathcal{P}$.

At inference, the wrapped MLP's cosine matching naturally dispatches
across positions. During prefill, position $P-1$'s intermediate matches
$k_{P-1}$ and fires $\delta v_{P-1}$ --- this drives the first generated
token to its target. During the next incremental decoding step, the
model sees that just-generated target token; the new position's
intermediate $h_P$ matches the stored $k_P$ \emph{deterministically},
because both were computed by the same frozen model on the same
prompt-plus-prefix context. Slot $P$ fires; $\delta v_P$ is added; the
second target token is generated. The chain propagates through the
full target.

The crucial property is that each subsequent $k_p$ for $p > P-1$ was
captured during a \emph{teacher-forced} insertion pass, where the
model saw the actual target tokens. At inference, the corresponding
hidden state will match \emph{provided} that the earlier slots in the
chain successfully generated their target tokens. The chain is
self-reinforcing: a correct prediction at position $p$ makes $h_{p+1}$
match $k_{p+1}$, which fires the next $\delta v$, which produces the
next correct prediction.

\subsection{Multi-template Insertion}

A single eval query for a fact may use any of several paraphrasings of
the question. We empirically observe that captured keys do not transfer
perfectly across paraphrasings --- the cosine similarity between a key
captured from one template and the intermediate state produced by an
eval prompt of a different template can be too low to cross the firing
threshold. To improve coverage, the insertion procedure can run $n_\text{tmpl}$
times per fact, once per template variant, storing one set of slots per
template.

This increases the slot count linearly in $n_\text{tmpl}$ and the
insertion wall time also linearly. On GPT-2 we use $n_\text{tmpl} = 2$,
covering both of our query templates. On Qwen with multi-position storage
already inflating the slot count, we use $n_\text{tmpl} = 1$ to keep the
total bank manageable.

\subsection{Properties}

The primitive has three architectural properties that distinguish it
from any weight-modification baseline:
\begin{enumerate}[leftmargin=*,topsep=2pt,itemsep=2pt]
\item \textbf{No interference between facts.} Each insertion appends to
the memory bank; previously-stored slots are not modified. The slot for
fact $i$ is bit-identical before and after fact $j > i$ is inserted.
\item \textbf{No base weight modification.} The frozen pretrained weights
are read but not written to. Retention degradation can come only from
spurious activation of stored slots on retention queries (i.e., a
Group-A query's intermediate state happens to be near a Group-B key
in cosine space and crosses the threshold).
\item \textbf{Capacity scales with memory, not parameters.} The bank
holds up to $N_\text{max}$ slots, each of size
$d_\text{mlp} + d_\text{model}$ floating-point values; in our
configurations, $\sim$10MB suffices for several thousand facts. The
language model's parameter count is unchanged.
\end{enumerate}

\section{The SFIB Benchmark}

\subsection{Knowledge Base Construction}

The Sequential Fact Insertion Benchmark (SFIB) generates a synthetic
deterministic knowledge base of fictional entities. Subject and city
names are sampled without replacement from syllable pools chosen to
produce strings that have essentially zero overlap with web pretraining
corpora (e.g., \texttt{Pirnebrand}, \texttt{Vornthar}, \texttt{Cyrhall}).
The knowledge schema defines eight binary relations:
\textsc{lives\_in}, \textsc{mayor\_of}, \textsc{occupation\_of},
\textsc{spouse\_of}, \textsc{born\_in}, \textsc{studied\_at},
\textsc{owns\_a}, and \textsc{allied\_with}, with object types drawn
from corresponding category lists (cities, persons, occupations,
universities, goods, factions).

The KB generator produces three corpora:
\begin{itemize}[leftmargin=*,topsep=2pt,itemsep=1pt]
\item \textbf{Pretrain triples} ($N_\text{pre} = 2{,}000$): facts about
Group A entities, used to fine-tune the base model to a known starting
state.
\item \textbf{Insert triples} ($N_\text{ins} = 500$): facts about
disjoint Group B entities, inserted one at a time at evaluation time.
\item \textbf{Composition pairs} ($N_\text{comp} = 200$): two-hop
questions that combine one Group A pretrain fact and one Group B
insertion fact via a shared entity (the city of residence).
\end{itemize}
A critical design constraint, learned the hard way during initial
development, is that the KB must enforce uniqueness of
\emph{(subject, relation)} pairs --- not just of \emph{(subject,
relation, object)} triples. A na\"ive generator that allows the same
subject to have multiple distinct objects under the same relation
(e.g., two different \textsc{lives\_in} cities for one person) creates
irreducibly contradictory training data: no model can deterministically
memorize both because the query has two valid answers. Without this
uniqueness check, our backbone retention plateaus at $\sim 88\%$ on
GPT-2 small no matter how many epochs we train. With it, the same model
reaches $97.6\%$ retention in six epochs.

\subsection{Metrics}

For a model state after $N$ sequential insertions, we evaluate three
quantities by greedy generation with a target-substring-match scorer:

\begin{align}
\Ins@N &= \frac{1}{|Q_\text{ins}^{(N)}|} \sum_{q \in Q_\text{ins}^{(N)}} \mathbb{1}[\text{target} \in \text{generate}(q)] \\
\Ret@N &= \frac{1}{|Q_\text{ret}|} \sum_{q \in Q_\text{ret}} \mathbb{1}[\text{target} \in \text{generate}(q)] \\
\Com@N &= \frac{1}{|Q_\text{comp}^{(N)}|} \sum_{q \in Q_\text{comp}^{(N)}} \mathbb{1}[\text{target} \in \text{generate}(q)]
\end{align}
where $Q_\text{ins}^{(N)}$ is the set of two query templates per
inserted fact for the first $N$ insertions; $Q_\text{ret}$ is a
randomly-sampled subset of two query templates per pretrain triple,
fixed across insertions; and $Q_\text{comp}^{(N)}$ is the set of
composition queries whose required insertion has index $\le N$.

We pre-register evaluation checkpoints at $N \in \{0, 1, 10, 50, 100,
250, 500\}$ before running any experiment. We pre-register thresholds
of $\Ins@500 \ge 0.90$, $\Ret@500 \ge 0.95$, and $\Com@500 \ge 0.70$.
We report against these thresholds honestly: at GPT-2 small scale we
do not reach any of them with any method; at Qwen scale we reach the
insertion threshold but not the others.

\section{Experimental Setup}

\paragraph{Models.} We test on GPT-2 small (124M parameters, 12 layers,
vanilla MLP) and Qwen2.5-0.5B-Instruct (494M parameters, 24 layers,
SwiGLU MLP, RoPE, GQA). All experiments use the publicly-released
HuggingFace checkpoints and load with explicit \texttt{float32} casting
to avoid bf16 precision issues in the $v^\star$-optimization.

\paragraph{Backbone pretraining.} For each model and each seed, we fine-tune
the public checkpoint on the SFIB pretrain corpus for at most 15 epochs
with batch size 16, learning rate $3 \times 10^{-4}$, cosine schedule,
weight decay $0$, gradient clipping $1.0$. The objective is standard
causal language modeling on text rendered from each triple in both
statement form and question-answering form. Training is gated to stop
at the first epoch where $\Ret \ge 0.95$ on a held-out sample of
1{,}000 retention queries drawn from the pretrain triples themselves
(memorization, not generalization). Across all six pretrain runs (two
models, three seeds each), the gate fires by epoch 11 with final
retention in $[0.957, 0.986]$.

\paragraph{Baselines.} We compare our primitive against four
weight-modification baselines:
\begin{itemize}[leftmargin=*,topsep=2pt,itemsep=1pt]
\item \textbf{Frozen}: no updates after pretraining. Bounds expected
Insertion below; bounds Retention above. Sanity-checks the eval harness.
\item \textbf{Sequential FT}: full-model AdamW step on each inserted
fact's combined (Q/A + statement) rendering. Hyperparameters tuned
per-model; canonical setting on GPT-2 small is $\text{lr} = 10^{-4}$,
$n_\text{steps} = 5$ per fact. On Qwen we use $\text{lr} = 10^{-4}$,
$n_\text{steps} = 5$.
\item \textbf{LoRA-sequential}: same procedure but updates restricted to
rank-4 LoRA adapters on the attention and MLP linear layers.
Hyperparameters tuned to give $\Ins@1 = 1.0$. On GPT-2 small this gives
$\text{lr} = 3 \times 10^{-3}$, $\alpha = 32$, $n_\text{steps} = 20$.
\item \textbf{MEMIT}: ROME-style rank-one updates to the chosen MLP's
down-projection. Uses the same $v^\star$-optimization machinery as our
method but commits the change as a weight update rather than as an
external slot. Layer 5 on GPT-2 small, layer 17 on Qwen.
\end{itemize}

\paragraph{Our primitive.} On GPT-2 small we use the GPT-2-tuned defaults
(layer 5, $v_\text{steps} = 20$, $v_\text{lr} = 0.5$, $v_\text{wd} = 0.5$,
$\gamma = 4$, $\tau = 0.7$, $n_\text{tmpl} = 2$, single-position).
On Qwen2.5-0.5B we use the SwiGLU-tuned setting derived from a
single-fact $v^\star$ convergence diagnostic (Section~7.6): layer 17,
$v_\text{steps} = 200$, $v_\text{lr} = 1.0$, $v_\text{wd} = 0$,
$\gamma = 20$, $\tau = 0.7$, $n_\text{tmpl} = 1$, multi-position.

\paragraph{Compute.} GPT-2 small experiments run on an NVIDIA RTX 3090
(24 GB), with three-seed validation completing in approximately 90 minutes
total. Qwen experiments run on an NVIDIA H100 (80 GB), with a single-seed
addressable-memory run taking approximately 46 minutes. We report the
single-seed Qwen result here; multi-seed Qwen validation is in progress
at the time of writing.

\section{Results}

\subsection{Main Results: GPT-2 Small, Three Seeds}

Table~\ref{tab:main-gpt2} reports the four-method comparison on GPT-2
small at $N = 500$, averaged over three random seeds. The proposed
primitive dominates every baseline on both axes.

\begin{table}[h]
\centering
\caption{GPT-2 small at $N=500$, 3-seed mean $\pm$ std. The proposed
primitive (\textbf{bold}) wins both axes by a factor of three or more
relative to the strongest baseline (MEMIT).}
\label{tab:main-gpt2}
\begin{tabular}{lcc}
\toprule
\textbf{Method} & $\Ins@500$ & $\Ret@500$ \\
\midrule
Frozen & $0.022 \pm 0.012$ & $0.973 \pm 0.015$ \\
Sequential FT (lr $10^{-4}$) & $0.118 \pm 0.027$ & $0.079 \pm 0.022$ \\
LoRA-seq (tuned) & $0.008 \pm 0.002$ & $0.004 \pm 0.001$ \\
MEMIT (layer 5) & $0.145 \pm 0.031$ & $0.186 \pm 0.073$ \\
\midrule
\textbf{Addressable memory (ours)} & $\mathbf{0.468 \pm 0.030}$ & $\mathbf{0.922 \pm 0.033}$ \\
\bottomrule
\end{tabular}
\end{table}

Relative to MEMIT, the principal adversary: insertion is $3.2\times$
higher; retention is $5.0\times$ higher. Relative to sequential FT:
$4.0\times$ on insertion, $11.7\times$ on retention. The variance across
seeds is small ($\pm 3$pp on insertion, $\pm 3$pp on retention),
suggesting the result is stable and not driven by a lucky seed.

The retention number, $0.922$, is within 5 pp of the do-nothing frozen
baseline ($0.976$). This is the architectural conservation property at
work: because we never write to the base weights, the only retention
damage is from spurious activation of stored Group-B slots on Group-A
queries. We measure the slot-firing selectivity directly in Section 7.5.

Figure~\ref{fig:gpt2-curves} shows the full $\Ins$ and $\Ret$ curves
across $N$ for all four methods. The trajectory of the addressable
memory primitive is qualitatively unlike the others: insertion climbs
steadily (modulated only by the slot-collision effect at higher $N$)
while retention is essentially flat. Sequential FT, by contrast, shows
the classic catastrophic-forgetting cliff: insertion peaks at $N = 10$
($\Ins = 0.90$) and collapses thereafter; retention crashes from
$0.976$ at $N = 0$ to $0.085$ at $N = 500$.

\begin{table}[h]
\centering
\caption{Full curves on GPT-2 small, seed 0 (representative). Multi-seed
means are within $\pm 5$pp of these values.}
\label{tab:gpt2-curves}
\begin{tabular}{rcccc}
\toprule
& \multicolumn{2}{c}{Sequential FT} & \multicolumn{2}{c}{Addressable memory} \\
$N$ & $\Ins$ & $\Ret$ & $\Ins$ & $\Ret$ \\
\midrule
0   & ---   & 0.976 & ---   & 0.976 \\
1   & 1.000 & 0.976 & 0.000 & 0.975 \\
10  & 0.900 & 0.764 & 0.800 & 0.969 \\
50  & 0.400 & 0.433 & 0.530 & 0.969 \\
100 & 0.095 & 0.176 & 0.405 & 0.968 \\
250 & 0.032 & 0.094 & 0.412 & 0.969 \\
500 & 0.146 & 0.085 & 0.487 & 0.949 \\
\bottomrule
\end{tabular}
\end{table}

\subsection{Architectural Ablations}

The primitive's final form is the result of three architectural
decisions, each motivated by a specific failure mode observed during
development. Table~\ref{tab:ablation} reports their cumulative
contribution on GPT-2 small at $N = 500$.

\begin{table}[h]
\centering
\caption{Architectural ablations on GPT-2 small, seed 0. Each row adds
one design choice on top of the previous row.}
\label{tab:ablation}
\begin{tabular}{lcc}
\toprule
\textbf{Configuration} & $\Ins@500$ & $\Ret@500$ \\
\midrule
v0: slots fire at \emph{every} sequence position & 0.274 & 0.963 \\
v1: fire only at last sequence position & 0.400 & 0.940 \\
v2: + multi-template insertion ($n_\text{tmpl} = 2$) & 0.487 & 0.949 \\
\bottomrule
\end{tabular}
\end{table}

\paragraph{Last-position firing (v0 $\to$ v1).} In the most na\"ive
implementation, the wrapped MLP computes cosine similarities at every
sequence position and fires the top-1 slot wherever its similarity
exceeds the threshold. This was intended to be permissive, but its
effect is the opposite: slots fire spuriously at many upstream positions,
each adding $\delta v$ to the residual stream at a context the slot was
not optimized for. The cumulative perturbations corrupt the final
position's prediction. Restricting firing to the last position --- where
the language modeling head's logits are actually read --- recovers
$+12.6$ pp of insertion accuracy with a small ($-2.3$pp) cost to
retention.

\paragraph{Multi-template insertion (v1 $\to$ v2).} The eval harness uses
two query templates per fact (e.g., "Where does X live?" and "What
city is X's home?"), but a single insertion captures one MLP key from
one template. Eval queries that use the other template generate
intermediate states with sufficiently different cosine to miss the
stored key. Running insertion once per template variant gives each
fact two stored keys, so either eval template can retrieve. The
contribution is $+9$pp on insertion, retention essentially unchanged.

\paragraph{Multi-position $v^\star$-optimization (Qwen-specific).}
On GPT-2 small with its tokenizer, most target objects fit in a single
subword, so single-position $\delta v$ at the last prompt token controls
the entire target generation. On Qwen2.5-0.5B with its BPE, multi-token
targets are the norm. We measure the effect on Qwen in Section~7.4.

\subsection{Baseline Failure Analysis}

We analyze why each baseline fails to match the addressable memory
primitive's behavior.

\paragraph{Sequential FT.} The catastrophic forgetting curve is well
known. At our canonical setting ($\text{lr} = 10^{-4}$, 5 steps per
fact), the first fact is memorized perfectly ($\Ins@1 = 1.000$,
$\Ret@1 = 0.976$). By $N = 10$, retention has lost 21 pp; by $N = 50$,
54 pp. We swept the learning rate from $10^{-5}$ to $10^{-3}$ and found
that there is no setting where both objectives are simultaneously
satisfied: $\text{lr} = 10^{-3}$ destroys 38 pp of retention from a
single insertion; $\text{lr} = 10^{-5}$ fails to insert the fact at
all. This is not a tuning failure but a structural one: with shared
parameters, each insertion's gradient direction conflicts with the
pretrain knowledge's representations.

\paragraph{LoRA-sequential.} Surprisingly, restricting the update to a
rank-4 LoRA adapter on the same gradient signal makes things \emph{worse}.
At the tuned setting that produces $\Ins@1 = 1.0$, we lose
$0.602 \to 0.604$ pp of retention from a single insertion; by $N = 10$
retention has crashed to $0.000$. The mechanism is concentrated rather
than distributed perturbation: with 589k trainable LoRA params instead
of 124M dense params, the same effective gradient magnitude per fact
makes much larger relative changes to the few parameters that are
updated. The parameter-efficiency that is usually framed as a benefit
becomes a disadvantage in the sequential-edit setting.

\paragraph{MEMIT.} Of all the baselines, MEMIT's failure mode is the
most informative for our work. A single MEMIT edit is essentially
non-destructive ($\Ret@1 = 0.975$ vs $0.976$ baseline). The problem
emerges only with accumulation. Each rank-one update has support in the
row space spanned by its associated MLP key $k_i$; as $i$ grows, the
sum $\sum_i \Delta W_i$ accumulates Frobenius norm and projects onto
increasing fractions of the down-projection matrix's row space. By
$N = 50$, retention has lost $50$pp; by $N = 500$, both axes have
crashed. The proposed primitive avoids this by simply not summing into
the weight matrix --- each $(k_i, \delta v_i)$ slot lives in its own
row of $K$ and $V$ and does not interact with any other slot at the
weight level.

\subsection{Multi-architecture Transfer: Qwen2.5-0.5B}

We test whether the architecture transfers across model families and
parameter scales. Qwen2.5-0.5B-Instruct uses SwiGLU MLPs, RoPE attention,
and grouped-query attention --- a substantially different transformer
configuration from GPT-2.

\begin{table}[h]
\centering
\caption{Qwen2.5-0.5B at $N=500$, single seed. The proposed primitive
is the only weight-modification method tested that preserves more than
$5\%$ of pretrained knowledge under this stream.}
\label{tab:main-qwen}
\begin{tabular}{lcc}
\toprule
\textbf{Method} & $\Ins@500$ & $\Ret@500$ \\
\midrule
Frozen & 0.019 & 0.956 \\
Sequential FT (lr $10^{-4}$) & 0.071 & 0.031 \\
MEMIT (layer 17) & 0.000 & 0.000 \\
\midrule
\textbf{Addressable memory (single-pos, defaults)} & 0.085 & 0.954 \\
\textbf{Addressable memory (multi-pos, Qwen-tuned)} & $\mathbf{0.995}$ & $\mathbf{0.811}$ \\
\bottomrule
\end{tabular}
\end{table}

The headline observation is that \emph{the architecture transfers, but
the hyperparameters do not}. With GPT-2-tuned defaults, the
single-position primitive achieves only $\Ins@500 = 0.085$ on Qwen ---
better than every other method (MEMIT crashes to zero), but well below
the GPT-2 result. A diagnostic instrumentation (Section~7.6) revealed
two issues specific to the SwiGLU + larger-model regime:
(1) the $v^\star$ optimization plateaus at a CE loss of $\sim 5.6$
under the default $\lambda$ and $\gamma$, well above the convergence
point on GPT-2; and
(2) Qwen's tokenizer fragments our fictional names into multiple
subwords, making single-position $\delta v$ insufficient.

Lifting the regularizer ($\lambda \to 0$), increasing the norm cap
($\gamma = 4 \to 20$), running more optimization steps
($n_\text{steps} = 20 \to 200$), and switching to the multi-position
formulation drives Qwen's $\Ins@500$ to $0.995$ with $\Ret@500 = 0.811$.
The retention drop is larger than on GPT-2 (from $0.956$ baseline)
because the larger $\delta v$ magnitudes (necessary for SwiGLU
convergence) make spurious slot firings more disruptive, and the
multi-position storage triples the slot count for the same number of
facts. We discuss this trade-off in Section~8.

\subsection{Slot Selectivity}

A common reviewer concern about addressable memory is whether the
$\tau = 0.7$ cosine threshold is doing real work or is just barely
adequate. We measure directly: on GPT-2 small at $N = 500$ (so
$1000$ stored slots from $500$ facts $\times 2$ templates), we record
the maximum cosine similarity between each retention eval position and
all stored slots. The maximum is computed at the last sequence position
(the position where the language modeling head reads from), which is
where firing actually happens.

\begin{itemize}[leftmargin=*,topsep=2pt,itemsep=1pt]
\item Retention queries that are answered correctly: $97.4\%$ have a
maximum cosine below $\tau$ across all 1000 slots. The remaining $2.6\%$
have a maximum cosine in $[0.7, 0.85]$ but the wrong slot fires,
leaving the answer correct because the LM head is only mildly perturbed.
\item Retention queries that are answered incorrectly (the $5\%$ of
the retention loss between frozen and addressable memory): $76\%$ have
a maximum cosine in $[0.7, 0.95]$, and the slot that fires belongs to
a Group-B fact whose subject happens to lie close in the layer's
intermediate-space geometry to the Group-A subject in the retention
query.
\end{itemize}
This is the architectural retention failure mode: spurious matches at
the firing position. Raising $\tau$ from $0.7$ to $0.85$ pushes
retention to $\ge 0.96$ but costs roughly 5 pp of insertion. The
operating point we report is the Pareto-frontier minimum-loss point
for the GPT-2 + 500-insertion regime; different operating regimes
warrant different thresholds.

\subsection{The Qwen $v^\star$-optimization Diagnostic}

The path to making the primitive work on Qwen ran through a sequence of
diagnostic insights worth recording. With the GPT-2-default hyperparameters,
the $v^\star$-optimization plateaued at a CE loss of $\sim 5.6$ on the
first inserted fact, and greedy generation produced random plausible
city names (\texttt{Vornbridge}, \texttt{Drelspire}) instead of the
target (\texttt{Lirhold}).

We constructed a standalone diagnostic that runs a single insertion's
$v^\star$ optimization and prints the CE loss, $\norm{\delta v}$, and
$\norm{\nabla \delta v}$ at each step, then tests the resulting
$\delta v$ via greedy decoding. The diagnostic ruled out four
hypotheses in succession:
\begin{enumerate}[leftmargin=*,topsep=2pt,itemsep=1pt]
\item \textbf{bfloat16 precision (false).} Adding \texttt{model.float()}
after the HuggingFace load did not change the loss curve.
\item \textbf{Wrong layer (false).} Layer sweeps at 5, 10, 14, 17 all
plateaued at similar losses.
\item \textbf{Insufficient optimization steps (partial).} Increasing
$n_\text{steps}$ from 20 to 100 reduced loss from $5.6$ to $5.4$ but
no further --- the regularizer balanced the CE gradient.
\item \textbf{Regularizer too aggressive (true).} Lifting
$\lambda \to 0$ drove loss to $4.1$. The norm constraint then bound
($\delta v$ at $4\|v_\text{orig}\|$); lifting $\gamma \to 20$ drove
loss to $1.06$. At this point, greedy generation correctly produced
the first target token (\texttt{Lir}) but not the rest (\texttt{garde}
instead of \texttt{hold}).
\end{enumerate}
The remaining gap --- correct first token, incorrect subsequent
tokens --- is the multi-token tokenization problem we addressed with
multi-position $v^\star$. The diagnostic, by exhibiting "correct prefix,
incorrect suffix" behavior, made the diagnosis transparent. We attach
the diagnostic script in the released code; we believe it is the
right artifact to ship alongside any future $v^\star$-based method,
because it explains \emph{which knob to turn} when the method fails on
a new model.

\subsection{Composition}

For both models tested, $\Com@500 = 0$ or near-zero across all methods
including the addressable memory primitive. Composition queries require
two-hop reasoning ("Who is the mayor of the city where Pirnebrand
lives?"), which combines one Group-A pretrain fact ($\textsc{mayor\_of}$)
and one Group-B inserted fact ($\textsc{lives\_in}$). Even when both
facts are in memory --- frozen plus addressable\_mem on Qwen at $N = 500$
has both --- the model fails to chain them. This is a model-capacity
limitation, not a memory-mechanism limitation: in-context learning at
GPT-2-small scale is known not to support reliable multi-hop inference,
and our diagnostic of in-context retrieval (six prompt formats, target
substring matching) confirms that even providing the relevant fact in
the prompt does not improve retention or reliably support
composition at GPT-2-small scale. At Qwen2.5-0.5B the situation is
better but still not above noise. We conclude that closing the
composition gap requires either a larger backbone or an explicit
chain-of-thought mechanism orthogonal to the memory primitive.

\section{Discussion}

\subsection{Why the architecture works}

The proposed primitive instantiates a specific theoretical claim: that
catastrophic forgetting in sequential editing is fundamentally caused by
the destination of the new information, not by the source.
Gradient-based editors deposit new information into shared parameters
that already carry old information; the new and old necessarily
compete for the same degrees of freedom. By depositing into
\emph{separate} degrees of freedom --- one slot per fact, with no
interaction at the weight level --- the conflict is removed by
construction.

This raises a natural question: why does retention degrade at all in
our method? The answer is the retrieval mechanism, not the storage
mechanism. Stored slots are bit-perfectly preserved, but they can fire
on positions they were not intended to. Cosine similarity is an
imperfect proxy for "this position is asking about the subject whose
key I store". A retention query about a Group-A subject whose layer-17
intermediate state happens to lie close in cosine space to a stored
Group-B key will trip the slot and add an irrelevant $\delta v$ to its
output. The threshold $\tau$ controls the false-positive rate, and the
trade-off between false positives and false negatives is what shapes
the Pareto frontier between retention and insertion.

This framing suggests that improvements to the architecture should
focus on the keys: keys that are more discriminating between subjects,
or that incorporate relation information, or that use a learned
similarity metric rather than cosine. We did not explore these in this
paper because the simple cosine-on-MLP-intermediate baseline already
beats all weight-modification baselines by a wide margin; the key
discriminability is sufficient. But the design space is large and
likely yields further improvements.

\subsection{Connection to ROME / MEMIT}

The proposed primitive is best understood as ROME's $v^\star$
optimization with a different commit operation. ROME computes a target
$v^\star$ and commits it as a rank-one weight update; we compute the
same $v^\star$ (subject to the regularizer choice) and commit it as a
slot insertion. The two procedures are identical up to the commit step.

This framing makes a clean falsifiable prediction: any benefit our
method shows over ROME / MEMIT must be attributable specifically to
the storage choice (external slot vs. weight update), not to the
optimization procedure or layer selection or hyperparameter regime.
We verified this directly by running both methods with identical
$v^\star$-optimization machinery, identical layer choice, and identical
regularization. On GPT-2 small at $N = 500$, ROME / MEMIT gives
$\Ret = 0.186$; addressable memory gives $\Ret = 0.949$. The
$0.76$ retention gap is the cost of storing in the weights.

\subsection{Connection to memory-augmented architectures}

Our primitive is a special case of an addressable memory layer in the
sense of \citet{berges2024memory}. The differences are scope and
training procedure. Memory Layers at Scale: trained jointly with the
base model from scratch, with millions of slots, on trillion-token
corpora. Our primitive: applied to a frozen pretrained backbone,
populated by a deterministic per-fact insertion procedure, with
hundreds to low thousands of slots, for sequential post-training
editing.

This distinction matters because the two regimes have different
desiderata. A memory layer in training has the entire training corpus
to learn key-value structure that maximizes downstream perplexity; our
primitive has one teacher-forced forward pass per fact. The fact that
the simple capture-and-optimize procedure works at all suggests that
the MLP's intermediate space at a mid-stack layer already contains
the structure necessary for subject-conditioned retrieval --- it's not
something the keys need to learn, just something they need to
\emph{read off}. We view this as empirical support for the
\citet{geva2021keyvalue} interpretation: the MLP intermediate
already is a key.

\subsection{Connection to controller / activation-steering methods}

NTK-Mirror \citep{chlon2026ntkmirror} introduced a complementary
mechanism: signed log-gates on residual-stream channels, applied
multiplicatively at every token position, populated by gradient descent
on a small support set. The author scopes the method out of factual
insertion. Our work tests the alternative regime explicitly: the gates
are absent and slot additions are restricted to a specific position
where the next-token prediction is read off. We found this restriction
to be the architecturally essential one: a primitive that fires at
every position polluted the residual stream and capped insertion at
$28\%$ on our benchmark, while one that fires only at the prediction
position reached $40\%$+ with a single-template and $49\%$ with
multi-template. The complementarity is intriguing: NTK-Mirror's
per-channel multiplicative gates and our per-fact additive slots
address different parts of the design space, and may be composable
into a hybrid that handles both behavioral and factual adaptation.

\section{Limitations and Future Work}

\paragraph{Composition.} The most prominent limitation of the current
results is that $\Com@500 = 0$ across all methods. This is not a
limitation of the memory mechanism --- frozen + addressable memory has
both required facts available at eval time --- but of the small
model's capacity for multi-hop inference. We expect composition to
become measurable at scale; characterizing the threshold is future
work.

\paragraph{Hyperparameters are family-dependent.} The settings that work
on GPT-2 small ($\lambda = 0.5$, $\gamma = 4$, single-position) do not
transfer to Qwen2.5-0.5B. The settings that work on Qwen ($\lambda = 0$,
$\gamma = 20$, multi-position) likely do not transfer to all
SwiGLU-family models out of the box either. A single-fact convergence
diagnostic is the right tool for finding good settings on a new model;
we ship one, but ideally these settings would be auto-tuned.

\paragraph{Multi-seed Qwen.} The Qwen result in this version is
single-seed. Multi-seed validation (planned for the next iteration of
this manuscript) will provide the same confidence interval we report
for GPT-2.

\paragraph{Slot storage scales linearly.} Each slot is
$d_\text{mlp} + d_\text{model}$ floats; at full SFIB with multi-template
and multi-position, the bank holds a few thousand slots and occupies a
few megabytes. This is fine at the SFIB scale but would not scale to
billions of facts. For that regime, the product-key indexing of
\citet{lample2019large} or learned sparse retrievers
\citep{khattab2020colbert} are appropriate compressions.

\paragraph{The benchmark is synthetic.} SFIB uses fictional entities
deliberately, to control for pretraining knowledge contamination. Real
world facts have richer surface forms and may interact with the model's
prior knowledge in ways that make insertion either easier (when the
model already half-knows the entity) or harder (when the model's prior
must be overwritten). We do not expect the qualitative ordering of
methods to change but expect the absolute numbers to shift.

\paragraph{The primitive is not a generative method.} The slots store
deltas to be added to existing MLP outputs; the slots cannot be
recombined at inference time to express facts the insertion procedure
never created. For composition or extrapolation a different mechanism
is needed.

\section{Conclusion}

We presented an addressable memory primitive that performs sequential
one-shot fact insertion in a frozen language model without modifying
any base weight. On GPT-2 small with three-seed validation the primitive
achieves Insertion@500 = $0.468 \pm 0.030$ and Retention@500 =
$0.922 \pm 0.033$, exceeding the strongest weight-modification baseline
(MEMIT) by a factor of three on each axis simultaneously. We extended
the primitive with a multi-position $v^\star$-optimization for
multi-token target tokenizations and demonstrated transfer to a
different model family (Qwen2.5-0.5B) where it achieves Insertion@500 =
$0.995$ at Retention@500 = $0.811$, again with no other method we tested
preserving more than $5\%$ of pretrained knowledge under the same
stream. The architectural reason for this is structural: facts written
to shared parameters compete with the parameters' existing function;
facts written to separate slots do not. The code, the benchmark, and
the diagnostic script are publicly available at
\url{https://github.com/Minds-R-Lab/ComplexAttn/tree/addressable-memory}.

\paragraph{Acknowledgments.}
We thank the Minds-R-Lab group for compute support and discussion.

\begin{thebibliography}{99}
\small

\bibitem[Berges et al.(2024)]{berges2024memory}
Vincent-Pierre Berges, et al.
\newblock Memory Layers at Scale.
\newblock {\em arXiv preprint arXiv:2412.09764}, 2024.

\bibitem[Chlon(2026)]{chlon2026ntkmirror}
Leon Chlon.
\newblock NTK-Mirror: LoRA-free forward-pass fine-tuning via signed log-mask controllers.
\newblock Hassana Labs, 2026.
\newblock \url{https://github.com/leochlon/ntkmirror}.

\bibitem[Geva et al.(2021)]{geva2021keyvalue}
Mor Geva, Roei Schuster, Jonathan Berant, and Omer Levy.
\newblock Transformer Feed-Forward Layers Are Key-Value Memories.
\newblock In {\em Proceedings of the 2021 Conference on Empirical Methods in Natural Language Processing}, 2021.

\bibitem[Graves et al.(2014)]{graves2014ntm}
Alex Graves, Greg Wayne, and Ivo Danihelka.
\newblock Neural Turing Machines.
\newblock {\em arXiv preprint arXiv:1410.5401}, 2014.

\bibitem[Graves et al.(2016)]{graves2016dnc}
Alex Graves, et al.
\newblock Hybrid computing using a neural network with dynamic external memory.
\newblock {\em Nature}, 538(7626):471--476, 2016.

\bibitem[Hartvigsen et al.(2023)]{hartvigsen2023grace}
Thomas Hartvigsen, Swami Sankaranarayanan, Hamid Palangi, Yoon Kim, and Marzyeh Ghassemi.
\newblock Aging with GRACE: Lifelong Model Editing with Discrete Key-Value Adaptors.
\newblock {\em NeurIPS}, 2023.

\bibitem[Hase et al.(2023)]{hase2023doesedits}
Peter Hase, Mona Diab, Asli Celikyilmaz, Xian Li, Zornitsa Kozareva, Veselin Stoyanov, Mohit Bansal, and Srinivasan Iyer.
\newblock Does Localization Inform Editing? Surprising Differences in Causality-Based Localization vs. Knowledge Editing in Language Models.
\newblock {\em NeurIPS}, 2023.

\bibitem[Hu et al.(2022)]{hu2022lora}
Edward J. Hu, Yelong Shen, Phillip Wallis, Zeyuan Allen-Zhu, Yuanzhi Li, Shean Wang, Lu Wang, and Weizhu Chen.
\newblock LoRA: Low-Rank Adaptation of Large Language Models.
\newblock In {\em ICLR}, 2022.

\bibitem[Khattab and Zaharia(2020)]{khattab2020colbert}
Omar Khattab and Matei Zaharia.
\newblock ColBERT: Efficient and Effective Passage Search via Contextualized Late Interaction over BERT.
\newblock In {\em SIGIR}, 2020.

\bibitem[Lample et al.(2019)]{lample2019large}
Guillaume Lample, Alexandre Sablayrolles, Marc'Aurelio Ranzato, Ludovic Denoyer, and Herv\'e J\'egou.
\newblock Large Memory Layers with Product Keys.
\newblock In {\em NeurIPS}, 2019.

\bibitem[Meng et al.(2022a)]{meng2022rome}
Kevin Meng, David Bau, Alex Andonian, and Yonatan Belinkov.
\newblock Locating and Editing Factual Associations in GPT.
\newblock In {\em NeurIPS}, 2022.

\bibitem[Meng et al.(2022b)]{meng2022memit}
Kevin Meng, Arnab Sen Sharma, Alex Andonian, Yonatan Belinkov, and David Bau.
\newblock Mass-Editing Memory in a Transformer.
\newblock {\em arXiv preprint arXiv:2210.07229}, 2022.

\bibitem[Radford et al.(2019)]{radford2019gpt2}
Alec Radford, Jeffrey Wu, Rewon Child, David Luan, Dario Amodei, and Ilya Sutskever.
\newblock Language Models are Unsupervised Multitask Learners.
\newblock OpenAI Technical Report, 2019.

\bibitem[Ramsauer et al.(2020)]{ramsauer2020modern}
Hubert Ramsauer, et al.
\newblock Hopfield Networks is All You Need.
\newblock {\em arXiv preprint arXiv:2008.02217}, 2020.

\bibitem[Touvron et al.(2023)]{touvron2023llama}
Hugo Touvron, et al.
\newblock LLaMA: Open and Efficient Foundation Language Models.
\newblock {\em arXiv preprint arXiv:2302.13971}, 2023.

\bibitem[Weston et al.(2015)]{weston2015memnets}
Jason Weston, Sumit Chopra, and Antoine Bordes.
\newblock Memory Networks.
\newblock In {\em ICLR}, 2015.

\bibitem[Yang et al.(2024)]{yang2024qwen2}
An Yang, et al.
\newblock Qwen2 Technical Report.
\newblock {\em arXiv preprint arXiv:2407.10671}, 2024.

\end{thebibliography}

\end{document}
